\documentclass[11pt,letterpaper]{article}
\usepackage[margin=1in]{geometry}
\usepackage[T1]{fontenc}
\usepackage{lmodern,amsmath,amssymb,amsthm,mathtools,array,longtable,booktabs}
\usepackage[hidelinks]{hyperref}
\usepackage{xurl}
\newtheorem{theorem}{Theorem}[section]
\newtheorem{lemma}[theorem]{Lemma}
\newtheorem{proposition}[theorem]{Proposition}
\newtheorem{corollary}[theorem]{Corollary}
\theoremstyle{definition}\newtheorem{definition}[theorem]{Definition}
\newcommand{\C}{\mathbb C}\newcommand{\Q}{\mathbb Q}\newcommand{\Z}{\mathbb Z}
\newcommand{\F}{\mathbb F}\newcommand{\tr}{\operatorname{tr}}
\newcommand{\rank}{\operatorname{rank}}\newcommand{\disc}{\operatorname{disc}}
\newcommand{\gb}[3]{\genfrac{[}{]}{0pt}{}{#1}{#2}_{#3}}
\newcommand{\res}[1]{[#1]_N}\newcommand{\eps}{\varepsilon}
\newcommand{\M}{\mathcal M}
\setlength{\emergencystretch}{3em}
\title{Small order Fourier synthesis by signed Weyl resolvent banks}
\author{}
\date{10 October 2026}
\pdfmapfile{+lm.map}
\pdfmapfile{+cm.map}
\pdfmapfile{+symbols.map}
\begin{document}
\maketitle
\begin{abstract}
We study additive representations of the negative Fourier matrix by phase-weighted averages of cyclic resolvents built from the four signed pairs of one fixed Weyl system. The counted object is a whole bank with one parameter pair and arbitrary phase weights. We give exact small-order lower and upper arguments, with all inverse guards and normalization conventions explicit. The accepted lower and upper results establish minima four, five, and eight at orders three, five, and seven, respectively, over the complex numbers. The order-seven construction lies in the splitting field of an explicit cubic over the seventh cyclotomic field. The method combines classical Gaussian-binomial expansion, finite-character projection, and moment recurrences with support restrictions and exact coefficient elimination. The certificates use characteristic-zero algebra or justified good reduction, not extrapolation from finite-field searches. The results concern representation count; they do not establish an arithmetic speedup for the Fourier transform.
\end{abstract}

\section{The synthesis problem and theorem status}

Fix an odd integer $N\geq3$, a field $k$ whose characteristic does not divide $N$, and a primitive $N$th root $q\in k$. On the basis $e_0,\ldots,e_{N-1}$ define
\[
S e_i=e_{i+1},\qquad T e_i=q^i e_i,\qquad W_{a,b}=S^aT^b,
\qquad F_N[k,t]=q^{-kt}.
\]
Indices are taken modulo $N$, and $[a]_N$ denotes the representative in $\{0,\ldots,N-1\}$. The Fourier matrix is unnormalized. Fractions in exponents of $q$, such as $1/4$, mean inverses in $\Z/N\Z$, never a choice of a complex square root.

\begin{definition}[Permitted bank]
For $(\eps,\delta)\in\{+1,-1\}^2$, put $\zeta=q^{\eps\delta}$. A bank is
\begin{equation}\label{eq:bank}
 B_{\eps,\delta}(\alpha,\beta;w)
 =\frac1N\sum_{j=0}^{N-1} w_j
 \bigl(I-\alpha S^\eps-\beta\zeta^j T^\delta\bigr)^{-1},
 \qquad \Delta=1-\alpha^N-\beta^N\ne0,
\end{equation}
with $\alpha,\beta,w_j\in k$. A representation is a finite sum of such banks. Each bank may have its own signs, parameters, and weights. One bank counts once even when some phase weights vanish. Zero banks may be omitted. Banks with identical signs and parameter pairs combine by adding their phase weights; different orientations may share parameters. No additional diagonal gauges, alternate Weyl generators, coordinate changes, products of banks, or free transformations are allowed.
\end{definition}

Write $\M_k(N;q)$ for the minimum number of banks in an exact representation of $F_N$ over $k$. In particular write $\M_{\C}(N;q)$ over the complex numbers. The results below hold for every primitive root of the indicated order; the values in the main statement do not depend on that choice.

\begin{theorem}[Accepted small-order assertions]\label{thm:main}
For every primitive root of the indicated order,
\[
 \M_{\C}(3;q)=4,\qquad \M_{\C}(5;q)=5,
 \qquad \M_{\C}(7;q)=8.
\]
The order-three equality holds over every field of characteristic not dividing three that contains $q$. The order-five equality holds over every characteristic-zero field containing $q$. The order-seven lower bound holds over every characteristic-zero extension field containing $q$; attainment holds over the explicit cubic splitting field in Section~\ref{sec:n7upper}, and hence over $\C$.
\end{theorem}

\paragraph{Source acceptance and field boundary.}
The order-seven equality combines two separate accepted receipts: the lower-eight obstruction and the final upper-eight literal construction. The latter accepts the frozen author witness over the splitting field $L$ of the specified cubic. No assertion that the witness lies in $\Q(q)$ is made. The accepted receipt hashes and the independent-review status of this assembled manuscript are recorded in \texttt{THEOREM\_STATUS.json}.

This is a synthesis of previously reviewed arguments, with the final assembled manuscript still subject to an independent review. The earlier four-bank classification and the general-size bounds are retained in Section~\ref{sec:general}; no exact general-$N$ formula follows from the three small cases.

\section{Classical identities and exact conventions}\label{sec:algebra}

For $0\leq\ell\leq n<N$ define
\[
 \gb n\ell\zeta=\prod_{i=1}^{\ell}
 \frac{1-\zeta^{n-\ell+i}}{1-\zeta^i},
 \qquad \gb n0\zeta=1.
\]
Every factor divided by is nonzero. If $A=\alpha S^\eps$ and $D=\beta\zeta^jT^\delta$, then $DA=\zeta AD$. The noncommutative Gaussian-binomial theorem gives
\[
 (A+D)^n=\sum_{\ell=0}^n\gb n\ell\zeta A^{n-\ell}D^\ell.
\]
At the primitive $N$th root, the mixed coefficients in degree $N$ vanish and
\[
 X_j^N=(\alpha^N+\beta^N)I,\qquad X_j=A+D.
\]
Consequently the legal inverse in~\eqref{eq:bank} is precisely
\begin{equation}\label{eq:inverse}
 (I-X_j)^{-1}=\Delta^{-1}\sum_{n=0}^{N-1}X_j^n.
\end{equation}
Multiplying either side by $I-X_j$ proves the assertion on both sides. This classical algebra is covered by the Potter--Sch\"utzenberger discussion in \cite[Theorems 1--2]{HMS} and the coefficient convention in \cite[Proposition 1]{Schlosser}.

Set
\[
 \gamma_\ell=\frac1N\sum_jw_j\zeta^{j\ell},\qquad
 \eta_\ell=\frac{\beta^\ell\gamma_\ell}{\Delta}.
\]
Substitution and finite-character orthogonality give
\begin{equation}\label{eq:expansion}
 B_{\eps,\delta}
 =\sum_{\ell=0}^{N-1}\eta_\ell
   \sum_{r=0}^{N-1-\ell}\alpha^r
   \gb{r+\ell}{\ell}{\zeta}W_{\eps r,\delta\ell}.
\end{equation}
In lower bounds we may allow all $\eta_\ell$ independently: this enlarges the legal class, including the possible $\beta=0$ stratum. We never assert that a legal bank with $\beta=0$ actually has arbitrary higher-sector amplitudes.

The Weyl matrices are a basis, since
\[
 \tr(W_{a,b}^{-1}W_{c,d})=N\,\mathbf1_{(a,b)=(c,d)}.
\]
Direct extraction of the Fourier coefficient gives
\begin{equation}\label{eq:fouriercoef}
 \phi_N(a,b)=\frac1N\sum_{t=0}^{N-1}q^{-bt-(t+a)t}
 = C_N q^{(a+b)^2/4},\qquad
 C_N=\frac1N\sum_tq^{-t^2}\ne0.
\end{equation}
Completing the square proves the second equality. To prove the last claim, put $G_\pm=\sum_tq^{\pm t^2}$. The invertible change $(j,l)\mapsto(l-j,l+j)$ modulo odd $N$ yields
$G_-G_+=\sum_{u,v}q^{uv}=N$. Thus $C_N$ is nonzero. These are elementary finite-character identities; a standard odd-dimensional Weyl and character framework is given in \cite[Section II A--D]{Gross}, with different Weyl phases.

For all small-order proofs below let $h_\ell=\eta_\ell/C_N$ denote a \emph{normalized} amplitude and $f_N(a,b)=q^{(a+b)^2/4}$ the normalized target. All constructions use $\beta=1$ and nonzero $\alpha$. Hence $\Delta=-\alpha^N\ne0$, and arbitrary specified $h_\ell$ are realized by the exact raw weights
\begin{equation}\label{eq:weights}
 w_j=C_N\Delta\sum_{\ell=0}^{N-1}h_\ell q^{-\eps\delta j\ell}.
\end{equation}
There is no additional factor $1/N$ in this formula: it is already in the definition of the bank.

\section{Support and the order-three minimum}\label{sec:n3}

\begin{lemma}[Four orientations are necessary]\label{lem:support}
Every Fourier representation in the permitted grammar contains each of the four orientations.
\end{lemma}
\begin{proof}
Equation~\eqref{eq:expansion} confines a bank to
\[
 A_{\eps,\delta}=\{(a,b):[\eps a]_N+[\delta b]_N\leq N-1\}.
\]
For interior coordinates $1\leq a,b\leq N-1$, the conditions for $++,+-,-+,--$ are, respectively,
\[
 a+b\leq N-1,\qquad a<b,\qquad b<a,\qquad a+b\geq N+1.
\]
In particular $(\eps t,\delta t)$ belongs exclusively to the indicated orientation for $1\leq t\leq(N-1)/2$. Its target coefficient is nonzero by~\eqref{eq:fouriercoef}. No number of banks of the other orientations can supply it. The four triangles do cover the square: the cross-sign triangles cover unequal interior coordinates, one same-sign triangle covers each equal pair by oddness, and the axes are covered directly. Coverage alone makes no assertion about coefficients.
\end{proof}

\begin{proposition}\label{prop:n3}
At order three four legal banks suffice, over every field of characteristic different from three containing a primitive cube root.
\end{proposition}
\begin{proof}
Take $\beta=1$ and $\alpha_{++}=1$, $\alpha_{+-}=q$, $\alpha_{-+}=q$, $\alpha_{--}=1$. All four guards equal $-1$. For the row $b=0$ define
\[
 c_k=\frac13\sum_{a=0}^2f_3(a,0)q^{-ka},\qquad k=0,1,-1.
\]
Set $h_{++,0}=c_0$, $h_{+-,0}=c_1$, $h_{-+,0}=c_{-1}$ and $h_{--,0}=0$. The corresponding coefficient sequences are $1,q^a,q^{-a},1$, so character orthogonality gives the row.

Next put
\[
 h_{\eps,\delta,1}
 =\frac{f_3(\eps,\delta)}{\alpha_{\eps,\delta}(1+q^{\eps\delta})}.
\]
Every denominator is nonzero and the four exclusive nonaxis coefficients are matched. For each $d=\pm1$ set
\[
 h_{+,-d,2}=f_3(0,d)-h_{+,d,1}-h_{-,d,1},
 \qquad h_{-,-d,2}=0.
\]
An $\ell=2$ sector has only $r=0$, so these assignments correct the two remaining axis entries without altering any earlier entry. Formula~\eqref{eq:weights} supplies actual legal phase weights. All nine Weyl coefficients agree, hence the original matrices agree. Lemma~\ref{lem:support} proves minimality.
\end{proof}

\section{The order-five lower bound}\label{sec:n5lower}

Assume four banks represent $F_5$. Lemma~\ref{lem:support} forces one bank of each orientation. Put
\[
 x=\alpha_{++},\quad y=\alpha_{-+},\quad
 u=\alpha_{--},\quad v=\alpha_{+-},\quad
 c=q+q^4,\quad d=1+q^4.
\]
All four nodes are nonzero because an exclusive coefficient must be supplied. Work with $f_5(a,b)=q^{4(a+b)^2}$. In $K=\Q(q)$, $c(c+1)=1$.

The two exclusive points of each orientation force, for its node $a$,
\begin{equation}\label{eq:n5forced}
 h_1=(1+q+q^3)/a,\qquad h_2=q^2/a^2.
\end{equation}
The nonaxis equations at $b=1$, $a=2,3$, then reduce to
\begin{equation}\label{eq:n5rows}
 x=cy^2-q^2,\qquad y=cx^2-q^2.
\end{equation}
Simultaneous reflection $(a,b)\mapsto(-a,-b)$ gives the same equations for $(u,v)$. Subtracting shows
$(x-y)(1+c(x+y))=0$. Each pair is therefore either diagonal, $x=y=t$ with
\[
 B(t)=ct^2-t-q^2=0,
\]
or nondiagonal with the fixed sum and product
\begin{equation}\label{eq:n5sp}
 s=x+y=-1/c=q^3+q^2,\qquad
 p=xy=1/c^2-q^2/c=1-q-2q^2-2q^3.
\end{equation}
For the latter, $s,p$ and $s^2-4p=-3+4q+7q^2+7q^3$ are nonzero.

At $b=2$, the outer coefficients $a=1,4$ determine
\[
 h_{+-,3}=-(1+d/x)/v,\qquad
 h_{--,3}=-(1+d/y)/u.
\]
Reflection determines the other two third-sector amplitudes. The actual coefficients at $(0,2)$ and $(0,3)$ now require
\begin{align}
 q^2(x^{-2}+y^{-2})-u^{-1}-v^{-1}
 -d\bigl((xv)^{-1}+(yu)^{-1}\bigr)&=q,\label{eq:n5axis1}\\
 q^2(u^{-2}+v^{-2})-x^{-1}-y^{-1}
 -d\bigl((xv)^{-1}+(yu)^{-1}\bigr)&=q.\label{eq:n5axis2}
\end{align}
No free sector can change these entries: $\ell=4$ affects only $b=\pm1$, and $\ell=0$ only $b=0$.

For a row-compatible pair define
\[
 A(x,y)=q^2(x^{-2}+y^{-2})+x^{-1}+y^{-1}.
\]
Using~\eqref{eq:n5rows} gives the exact identity
\begin{equation}\label{eq:n5mixed}
 A(x,y)=2c-\frac{(x+y)(x-y)^2}{x^2y^2}.
\end{equation}
Subtracting the axis equations forces $A(x,y)=A(u,v)$. A diagonal pair has $A=2c$, while a nondiagonal pair does not. Mixed types are impossible.

\paragraph{Two diagonal pairs.}
Write $x=y=t$, $u=v=w$ and set $a=2c-q$. Equation~\eqref{eq:n5axis1}, using $B(t)=0$, is equivalent to
\[
 D(t,w)=a t w-2(t+w)-2d=0.
\]
Indeed the axis residual times $tw$ is $D(t,w)-2(w/t)B(t)$. If $t\ne w$, their sum is $1/c$ and product $-q^2/c$, whence
\begin{equation}\label{eq:n5distinct}
 cD(t,w)=-2-2q+2q^2-q^3\ne0.
\end{equation}
If $t=w$, put
\[
 l_1=2q^3+2q^2-q+2,\qquad l_0=q^3+2q^2+2q,
 \qquad L(t)=cD(t,t)-aB(t)=l_1t+l_0.
\]
The polynomial identity
\begin{equation}\label{eq:n5bezout}
 l_1^2B(t)-(cl_1t-cl_0-l_1)L(t)
 =1+4q-11q^2+6q^3\ne0
\end{equation}
excludes a common zero without dividing by $l_1$.

\paragraph{Two nondiagonal pairs.}
Both have the same $s,p$ in~\eqref{eq:n5sp}; their unordered entries agree. Thus $(u,v)$ equals either $(x,y)$ or $(y,x)$. Multiplying the residual of~\eqref{eq:n5axis1} by $p^2$ gives respectively
\begin{align}
 q^2(s^2-2p)-sp-2dp-qp^2&=-5-5q+5q^3\ne0,\label{eq:n5non1}\\
 q^2(s^2-2p)-sp-d(s^2-2p)-qp^2&=-9-19q-17q^2-5q^3\ne0.\label{eq:n5non2}
\end{align}
Every stated nonzero residue is a nonzero rational polynomial in $q$ of degree below four. Since $\Phi_5(X)=1+X+\cdots+X^4$ is irreducible over $\Q$ (shift by one and use Eisenstein at five), none can vanish. The cases exhaust every complex choice of the nodes, not merely nodes in $K$. Therefore four banks are impossible over any characteristic-zero extension containing $q$.

\section{The order-five upper bound}\label{sec:n5upper}

We now specify five legal banks without requiring a search. All have $\beta=1$ and, in order, their orientations and nodes are
\begin{equation}\label{eq:n5nodes}
 (++,1),\ (++,z),\ (-+,q),\ (+-,t),\ (--,t),
 \quad z=-2-3q-3q^2-q^3,\quad t=1+q.
\end{equation}
Use $H=q/(1+q)=1+q+q^3$ and $c=q+q^4$.

Here is a complete exact amplitude recipe, which also explains why the nodes work. Temporarily write the nodes as $(x,z,y,v,u)$ in these five positions. Define
\[
 M(y)=cy^2-q^2,\qquad N(y)=(y+q^2)/c.
\]
For $b=1$ the singleton $-+$ amplitude is $H/y$. The two forward amplitudes can be taken as
\[
 h_{x,1}=\frac{H(z-M(y))}{x(z-x)},\qquad
 h_{z,1}=\frac{H(M(y)-x)}{z(z-x)}.
\]
They satisfy the three nonaxis moment equations exactly when
\begin{equation}\label{eq:n5condition}
 N(y)=(x+z)M(y)-xz.
\end{equation}
For $b=4$, the $+-$ and $--$ first-sector amplitudes are $H/v,H/u$; their nonaxis conditions are
\begin{equation}\label{eq:n5negative}
 u=cv^2-q^2,\qquad v=cu^2-q^2.
\end{equation}
These statements follow directly by substituting~\eqref{eq:expansion}; alternatively the three moments in the first row are $H,HM(y),HN(y)$ and obey the two-node recurrence.

For each edge row the opposite-delta $\ell=4$ sectors affect only $a=0$. At $b=1$ put the $+-$ fourth-sector amplitude equal to $f_5(0,1)$ minus the sum of the three known first-sector amplitudes, and set the $--$ fourth-sector amplitude to zero. At $b=4$ put the first $++$ fourth-sector amplitude equal to $f_5(0,4)-H/v-H/u$, and set the other two positive-delta fourth-sector amplitudes to zero. These choices match the axes.

It remains to solve the three rows $b=0,2,3$. With rows indexed by $a=0,\ldots,4$ and bank columns in the stated order, set $G=1+q+q^2$ and $P=1+q^4+q^3$, and use
\[
 M_0=\begin{pmatrix}
1&1&1&1&1\\x&z&y^4&v&u^4\\x^2&z^2&y^3&v^2&u^3\\
x^3&z^3&y^2&v^3&u^2\\x^4&z^4&y&v^4&u
\end{pmatrix},
\]
\[
 M_2=\begin{pmatrix}
1&1&1&1&1\\Gx&Gz&0&-qv&0\\q^2x^2&q^2z^2&0&0&0\\
0&0&q^3y^2&0&0\\0&0&Py&0&-q^4u
\end{pmatrix},\quad
 M_3=\begin{pmatrix}
1&1&1&1&1\\-q^4x&-q^4z&0&Pv&0\\0&0&0&q^3v^2&0\\
0&0&0&0&q^2u^2\\0&0&-qy&0&Gu
\end{pmatrix}.
\]
For each of these three rows define its amplitude vector as $M_b^{-1}(f_5(a,b))_{a=0}^4$, assigning in each bank the sector $\ell=[\delta b]_5$. Different rows use different sectors, so these assignments cannot interfere.

Specialize $x=1$, $y=q$, $u=v=t=1+q$ and $z$ as in~\eqref{eq:n5nodes}. Direct reduction modulo $\Phi_5$ gives
\[
 M(q)=q-q^2+q^3,\quad N(q)=1+2q+2q^2+q^3,
 \quad z=\frac{N(q)-M(q)}{M(q)-1},\quad ct^2-t-q^2=0.
\]
Thus~\eqref{eq:n5condition} and~\eqref{eq:n5negative} hold. The nodes are nonzero, $z\ne1$, and
\begin{align*}
 \det M_0&=-1093-5351q-6894q^2-3592q^3,\\
 \det M_2&=31+21q-17q^2-30q^3,\\
 \det M_3&=6+19q+21q^2+9q^3.
\end{align*}
They are nonzero by irreducibility. This proves that the recipe gives every one of the 25 normalized amplitudes in $K$. The guards in bank order are
\[
 -1,\quad 269+760q+795q^2+325q^3,\quad -1,
 \quad3-5q^2-5q^3,\quad3-5q^2-5q^3,
\]
all nonzero, since they equal $-\alpha^5$. Formula~\eqref{eq:weights} now specifies all 25 raw phase weights. Hence the five banks give $F_5$. Together with Section~\ref{sec:n5lower}, this proves the order-five part of Theorem~\ref{thm:main}.

The frozen literal witness records the same construction with a fully expanded 25-entry amplitude table and raw weight table. Independent verification used exact elimination on every original $5\times5$ branch for all four primitive-root conjugates, not just the coefficient formula. The recipe above is itself complete and does not require executing that witness generator.

\section{The order-seven lower bound}\label{sec:n7lower}

We prove that every delta sign requires at least four banks, regardless of the number on the other side. All amplitudes remain independently free in this impossibility argument. Work over $K=\Q(q)$ with $q$ primitive of order seven and $f_7(a,b)=q^{2(a+b)^2}$.

First suppose exactly three banks have positive delta, with nodes $x,z$ in orientation $++$ and $y$ in orientation $-+$. At $b=1$ all negative-delta banks have $\ell=6$ and vanish for $a=1,\ldots,6$. Put $H=q/(1+q)$. The endpoint $a=6$ forces $h_y y=H$, hence $y\ne0$. At $a=1$ the equation is $h_xx+h_zz=H$. Set
\[
 A=h_xx/H,\qquad B=h_zz/H,\qquad A+B=1.
\]
No division by $x$ or $z$ is used. For $a=2,3,4,5$ set
\[
 m_{a-1}(y)=\frac{q^{2(a+1)^2}}{H[a+1]_q}
   -\frac{[8-a]_{q^{-1}}}{[a+1]_q}y^{6-a},
 \qquad [j]_q=1+q+\cdots+q^{j-1},\quad m_0=1.
\]
The nonaxis equations say $m_j=Ax^j+Bz^j$ for $j=1,2,3,4$. Explicitly,
\begin{align*}
 m_1&=-1-q-q^4-q^5-(q^2+q^5)y^4,\\
 m_2&=1+q+q^3+q^4-(1+q^3+q^4)y^3,\\
 m_3&=-q^2-(1+q^2+q^3+q^4+q^5)y^2,\\
 m_4&=-q^2-q^3-(q^2+q^3+q^4+q^5)y.
\end{align*}
Define
\begin{equation}\label{eq:n7DF}
 D=m_2-m_1^2,\quad S_n=m_3-m_1m_2,\quad
 P_n=m_1m_3-m_2^2,\quad R=Dm_4-S_nm_3+P_nm_2.
\end{equation}
For $s=x+z$, $p=xz$, the two-atom recurrence implies the polynomial identities
\begin{equation}\label{eq:n7rec}
 S_n=sD,\qquad P_n=pD,\qquad R=0.
\end{equation}
They hold even when $D=0$, $x=z$, either forward node or amplitude vanishes, or both atoms coincide. No genericity hypothesis is being imposed.

A second isolated row uses $b=2$. Every negative-delta bank has $\ell=5$ and vanishes at $a=2,3,4,5$, independently of its parameters and how many such banks occur. If $k_x,k_z,k_y$ are the new second-sector amplitudes, the endpoint $a=5$ forces
\[
 k_y=\frac1{y^2\gb42{q^{-1}}}.
\]
The remaining three equations become $k_xx^j+k_zz^j=n_j(y)$ for $j=2,3,4$, where
\begin{align*}
 n_2&=q^4/\gb42q=2+q+q^3+2q^4,\\
 n_3&=\frac{q-y^2\gb62{q^{-1}}/\gb42{q^{-1}}}{\gb52q}
      =-1-q^3+(-3-2q+q^2-2q^3-3q^4)y^2,\\
 n_4&=\frac{q^2-y\gb52{q^{-1}}/\gb42{q^{-1}}}{\gb62q}
      =q^5-(1+q+q^3+q^4)y.
\end{align*}
Thus $n_4-sn_3+pn_2=0$. Multiplication by $D$, using~\eqref{eq:n7rec}, forces a second polynomial condition
\begin{equation}\label{eq:n7F}
 E(y)=Dn_4-S_nn_3+P_nn_2=0.
\end{equation}
Here $R$ and $E$ have degree nine. (The source certificates name $E$ as $F$; the renaming avoids confusion with the Fourier matrix.)

\subsection{A complete short certificate of incompatibility}

Specialize $q\mapsto7$ in $\F_{29}$. This element has order seven. Every denominator in the original row formulas has nonzero image. For example $H\mapsto19$, and $\gb42q,\gb52q,\gb62q,\gb42{q^{-1}}$ map respectively to $8,9,23,18$. In increasing powers of $Y$, direct reduction of the displayed moment formulas gives
\begin{align*}
 \overline R&=(23,26,9,23,13,0,11,16,10,8),\\
 \overline E&=(28,5,17,4,2,15,24,10,28,3).
\end{align*}
The coefficient vectors
\[
 U=(21,24,2,5,25,23,15,2,18),\qquad
 V=(18,28,0,22,8,2,16,25,10)
\]
satisfy the literal polynomial identity
\begin{equation}\label{eq:n7modbez}
 U\overline R+V\overline E=1\quad\hbox{in }\F_{29}[Y].
\end{equation}
This identity can be verified by ordinary multiplication of the four displayed finite lists. The leading coefficients of both polynomials survive. In particular, multiplication by $\overline E$ in the rank-nine quotient by the monic $\overline R$ is invertible (its independently recorded determinant is $19$).

For completeness, the valid characteristic-zero deduction is as follows. Let $\mathcal O=\Z[q]_{(29,q-7)}$. The row-formula denominators and leading coefficient of $R$ are units of this local ring. Hence
\[
 \mathcal A=\mathcal O[Y]/(R_{\rm monic})
\]
is finite free of rank nine, and multiplication by $E$ has a matrix over $\mathcal O$. Its reduction is invertible by~\eqref{eq:n7modbez}; its determinant is therefore a unit in $\mathcal O$ and in particular nonzero in $K$. Thus $E$ is a unit in $K[Y]/(R)$. Evaluation at a common root of $R$ and $E$ in any extension field would send a unit to zero, a contradiction. Squarefreeness of $R$ is unnecessary. This is good reduction of a nonzero norm, not an inference from absence of finite-field points.

The independent certificate additionally contains all rational coefficients of an exact identity $U_0R+V_0E=1$ in $K[Y]$, checked by multiplication. It is an alternative to, not an additional assumption for, the short good-reduction proof above.

\subsection{Every multiplicity and degenerate case}

Transposition sends $(\eps,\delta)$ to $(-\eps,\delta)$ and reindexes $j$ to $-j$, preserving parameters, guards, and the symmetric target $F_7$. It converts the positive-delta split $(1,2)$ to the excluded split $(2,1)$. If $P e_i=e_{-i}$, conjugation by $P$ sends signs to $(-\eps,-\delta)$, preserves the phase convention, and satisfies $PF_7P=F_7$. It transfers the obstruction to either split with negative delta.

Every orientation is necessary by Lemma~\ref{lem:support}. A side with only two banks therefore has one of each epsilon sign and can be padded to the excluded three-bank case by a legal zero bank: $\alpha=\beta=1$, $w_j=0$, $\Delta=-1$. A side with fewer than two banks already lacks an orientation. Consequently each side requires at least four banks, proving the lower bound eight.

This exhaustion permits arbitrary opposite-side parameters, zero or repeated forward nodes, zero forward amplitudes, and $D=0$. The singleton $y=0$ is excluded by its nonzero endpoint; the sole division by $y^2$ is therefore justified. No node-difference, discriminant, or unproved rank assumption enters the lower bound.

\section{The exact order-seven upper construction}\label{sec:n7upper}

\paragraph{Accepted scope.}
The final independent upper-bound receipt accepts the frozen construction below, including all guards, all 56 original phase resolvents, and all 49 original Fourier entries. Combining it with Section~\ref{sec:n7lower} proves the order-seven equality in Theorem~\ref{thm:main}.

Let $K=\Q(q)$ and define
\begin{align*}
 s_2&=(-8+35q+25q^2+28q^3+q^4-17q^5)/29,\\
 s_3&=(-3-5q-16q^2-4q^3+4q^4+19q^5)/29,\\
 g(t)&=t^3+s_2t-s_3.
\end{align*}
Let $L$ be its splitting field over $K$, and let $x_1,x_2,x_3$ be its roots in any order. With $\beta=1$ throughout, choose
\begin{equation}\label{eq:n7banks}
 (++,x_1),(++,x_2),(++,x_3),(-+,1),
 (+-,1),(--,x_1),(--,x_2),(--,x_3).
\end{equation}
The node product $s_3$ is nonzero, and
\[
 \disc(g)=\frac{665859+635064q+863499q^2+425987q^3+275146q^4-15036q^5}{24389}\ne0.
\]
All these nonzero claims follow from $\deg\Phi_7=6$ and the displayed nonzero numerators of degree below six. Hence the roots are distinct and nonzero, $[L:K]\leq6$, and the six root guards $-x_i^7$ and two singleton guards $-1$ are legal. There is no assertion that $L=K$.

\subsection{Exact edge construction}

Set $H=q/(1+q)$ and use the order-seven moments from Section~\ref{sec:n7lower} with $y=1$. They are
\begin{align*}
 m_0&=1,& m_1&=-1-q-q^2-q^4-2q^5,& m_2&=q,\\
 m_3&=-1-2q^2-q^3-q^4-q^5,&
 m_4&=-2q^2-2q^3-q^4-q^5.
\end{align*}
Now $D=m_2-m_1^2=1+4q-q^2+3q^4\ne0$. The displayed cubic coefficients satisfy
\[
 s_2=(m_1m_3-m_4)/D,\qquad s_3=m_3+s_2m_1,
\]
and therefore
\[
 m_3=-s_2m_1+s_3m_0,\qquad m_4=-s_2m_2+s_3m_1.
\]
The invertible Vandermonde on the three distinct roots gives unique weights $A_i$ with $\sum_iA_ix_i^j=m_j$ for $j=0,1,2$. The cubic recurrence supplies $j=3,4$. Thus $h_i=HA_i/x_i$ and singleton amplitude $H$ match every nonaxis entry at $b=1$.

The first three moments of these actual normalized amplitudes are
\[
 K_0=H(m_2+s_2)/s_3,\qquad K_1=H,\qquad K_2=Hm_1,
\]
because $1/x_i=(x_i^2+s_2)/s_3$. Assign the opposite-delta singleton's sixth-sector amplitude $q^2-K_0-H$ and the opposite triple's sixth-sector amplitudes zero. This supplies the axis. Simultaneous reflection supplies the entire row $b=6$.

\subsection{The other rows and their exact nonzero minors}

Let $p_r=(p_{r0},p_{r1},p_{r2})$ be the coefficient vector of $t^r$ modulo $g$ in the basis $1,t,t^2$. Thus
\[
 p_0=(1,0,0),\quad p_1=(0,1,0),\quad p_2=(0,0,1),\qquad
 p_r=-s_2p_{r-2}+s_3p_{r-3}\quad(r\geq3).
\]
For each orientation and $(a,b)$ let
\[
 r=[\eps a]_7,\quad \ell=[\delta b]_7,\quad
 G_{\eps,\delta}(a,b)=
 \begin{cases}\gb{r+\ell}{\ell}{q^{\eps\delta}},&r+\ell\leq6,\\0,&r+\ell>6.\end{cases}
\]
For fixed $b$, define the $7\times8$ matrix $M_b$ whose row $a$ is
\[
 \bigl(G_{++}p_{a0},G_{++}p_{a1},G_{++}p_{a2},
 G_{-+},G_{+-},
 G_{--}p_{[-a]_7,0},G_{--}p_{[-a]_7,1},G_{--}p_{[-a]_7,2}\bigr).
\]
Each $G$ uses its own orientation and its own $r,\ell$. Its unknown vector is
\[
 z_b=(K_{++,0},K_{++,1},K_{++,2},h_{-+},h_{+-},
 K_{--,0},K_{--,1},K_{--,2}).
\]
This is precisely the original coefficient system with each triple replaced by its three power moments; the two singleton nodes are one.

For $b=0,2,3$, retain respectively the following zero-based columns, set the omitted coordinate to zero, and solve the resulting square system against $f_b[a]=q^{2(a+b)^2}$:
\[
 I_0=(0,1,2,3,5,6,7),\quad
 I_2=(0,1,2,3,4,5,6),\quad
 I_3=(0,1,2,3,4,6,7).
\]
In increasing row order their determinants are
\begin{align*}
 d_0={}&-\bigl(370501134818472+889546502909192q
 +1166936870814012q^2\\
 &\hspace{15mm}+994193630954024q^3+501377982443436q^4
 +58844381586428q^5\bigr)/500246412961,\\
 d_2={}&(726-4387q-12368q^2-16345q^3-13989q^4-6744q^5)/841,\\
 d_3={}&1+3q+5q^2+5q^3+3q^4+q^5.
\end{align*}
They are nonzero by irreducibility of $\Phi_7$, so these three exact solves exist in $K$.

Define $J(z_0,z_1,z_2,z_3,z_4,z_5,z_6,z_7)
=(z_5,z_6,z_7,z_4,z_3,z_0,z_1,z_2)$. Put $z_4=Jz_3$, $z_5=Jz_2$, and
\[
 z_1=(K_0,H,Hm_1,H,q^2-K_0-H,0,0,0),\qquad z_6=Jz_1.
\]
The reflection identity proves the newly assigned rows, and the edge argument proves $b=1,6$. Hence all seven systems $M_bz_b=f_b$ are solved. The exact row ranks are $(7,5,7,7,7,7,5)$; the two deficient edge rows are handled by explicit compatibility, not by inferring solvability from their ranks.

\subsection{All sector amplitudes and original phase weights}

For a triple of moments $k=(k_0,k_1,k_2)$ define at a root $x_i$
\begin{equation}\label{eq:n7amplitude}
 A_k(x_i)=\frac{k_2+x_ik_1+(x_i^2+s_2)k_0}{g'(x_i)}.
\end{equation}
The denominator is nonzero by separability. Lagrange interpolation implies
$\sum_iA_k(x_i)x_i^d=k_d$ for $d=0,1,2$: one may verify it from
$\sum_i1/g'(x_i)=\sum_ix_i/g'(x_i)=0$ and $\sum_ix_i^2/g'(x_i)=1$, together with the cubic relation. Thus in row $b$ assign sector $\ell=[\delta b]_7$ by
\[
 h_{++,i,\ell}=A_{(z_{b0},z_{b1},z_{b2})}(x_i),\quad
 h_{-+,\ell}=z_{b3},\quad h_{+-,\ell}=z_{b4},\quad
 h_{--,i,\ell}=A_{(z_{b5},z_{b6},z_{b7})}(x_i).
\]
For each bank the map $b\mapsto[\delta b]_7$ is a permutation, so this defines all sectors once. Formula~\eqref{eq:weights} gives all 56 original phase weights. Root permutations only permute identically oriented banks. This is a complete finite algebraic specification, with no numerical root selection and no omitted amplitudes.

All divisions are legitimate: the field characteristic is zero; the scalar Gaussian factors, $1+q$, $D$, $s_3$, and $d_0,d_2,d_3$ are nonzero; distinct roots give nonzero $g'(x_i)$ and Vandermonde denominators; and the branch guards were checked above. The denominators $29$ in the expanded cubic coefficients must not be reduced termwise modulo $29$. No finite-field point is being lifted for this existence argument.

\subsection{A literal inverse certificate}

The certificate also verifies the original matrix equality independently of merely substituting into a coefficient table. Work in the separable cubic algebra $E=K[t]/(g)$. For $a_0+a_1t+a_2t^2\in E$,
\[
 \operatorname{Tr}_{E/K}(a_0+a_1t+a_2t^2)=3a_0-2s_2a_2.
\]
After passing to $L$, this trace is the sum over the three root evaluations, even if $g$ is reducible over $K$. Thus it sums a triple of banks exactly.

For every phase and each triple orientation construct the literal $7\times7$ matrix $X=tS^\eps+q^{\eps\delta j}T^\delta$ over $E$. Check $X^7=(t^7+1)I$ and both products of $I-X$ with $(I+\cdots+X^6)/(-t^7)$. Do the two singleton orientations at $t=1$. Four groups times seven phases give 28 algebra checks representing all 56 inverses: 42 in the triples and 14 singleton branches. Multiplying by the exact weights, taking traces for the triples, and adding gives all 49 entries $q^{-kt}$. The literal witness file records reduced sector polynomials, guard inverses, transformed amplitudes, and all exact checks. These data are a reproducible control of the proof, not a replacement for a legal-inverse argument.

\section{What the small cases do and do not settle}\label{sec:general}

The accepted earlier results imply, for any odd $N\geq3$,
\begin{equation}\label{eq:generalbound}
 \max\{4,2\lceil(N-1)/4\rceil\}\leq\M_{\C}(N;q)\leq4N.
\end{equation}
The lower mechanism is short enough to record. On $b=1$, nonaxis entries receive only positive-delta banks. A forward bank contributes
\[
 \eta_1\alpha^a\frac{1-q^{a+1}}{1-q},
\]
a combination of $\alpha^a$ and $(q\alpha)^a$. A backward bank with nonzero $\alpha$ contributes
\[
 \eta_1\alpha^{N-a}\frac{1-q^{a-1}}{1-q^{-1}},
\]
a combination of $(\alpha^{-1})^a$ and $(q\alpha^{-1})^a$. Zero nodes contribute zero on the whole interval. A sum of $r_+$ such banks therefore has Hankel rank at most $2r_+$. The Fourier Hankel matrix $H_{ij}=\phi_N(i+j+1,1)$, $0\leq i,j<(N-1)/2$, factors into nonzero diagonal matrices and a Vandermonde with nodes $q^{i/2}$. It has full rank $(N-1)/2$. Reflection treats negative delta. This proves the nonconstant lower bound. The exponential recurrence and Hankel factorizations are classical Prony ingredients \cite[Section 2]{PottsTasche}\cite[Theorem 2.1]{PeterPlonka}.

For the upper bound assign each Weyl coefficient to one containing triangle. In local coordinates divide the assigned value by $\gb{r+\ell}{\ell}\zeta$, setting unassigned entries to zero, to obtain $t_{r,\ell}$. Use $N$ nodes $\alpha_j=q^j$, $\beta=1$ per orientation, all with guard $-1$, and set
\[
 \eta_{j,\ell}=\frac1N\sum_{r=0}^{N-1}t_{r,\ell}q^{-jr}.
\]
Character orthogonality matches the assigned triangle; arbitrary amplitudes are legal by the weight conversion. Four orientations give $4N$ banks for any target matrix. Thus the order of growth of Fourier bank count is already $\Theta(N)$, but the exact count remains unsettled in general.

The earlier order-nine result strengthens the lower bound there to six. For completeness, if just two banks have positive delta, support puts one in each epsilon orientation. Their four row bases are $x,qx,z,qz$. Write $A=x+z$, $B=xz$, and $P(v)=\prod_{u\in\{x,qx,z,qz\}}(v-u)=\sum_{k=0}^4p_kv^k$. Then
\[
 p_3=-(1+q)A,\quad p_2=qA^2+(1+q^2)B,
 \quad p_1=qBp_3,\quad p_0=q^2B^2.
\]
The four valid Fourier row recurrences determine these coefficients, modulo $\Phi_9(q)=q^6+q^3+1$, as
\[
 p_3=-q^5+q^3-q^2-q+1,\quad p_2=-q^4+2q^2-1,
 \quad p_1=q^4-q^2+q+1=qp_3,\quad p_0=q^2.
\]
One has $p_3(q^2+1)=1$, so $B=1$. The $p_2,p_3$ equations require
\[
 E=(1+q)^2(p_2-1-q^2)-qp_3^2=q^5-q^4+2q^2-2q-3=0,
\]
whereas $E(2q^5-q^4+3q^2-3q+2)=3$ modulo $\Phi_9$, a contradiction in characteristic zero. Reflection gives three banks on the other side as well. This does not prove six-bank attainment.

Together with the small-order exclusions, the earlier Hankel bound for odd $N\geq11$ gives the preserved four-bank classification: for odd $N\geq3$ in characteristic zero, four banks suffice exactly at $N=3$. The order-five and order-seven equalities refine that classification without altering earlier cases. Exact minima for $N\geq9$, an order-seven witness wholly over $\Q(q)$, and positive-characteristic counterparts of the characteristic-zero small-order arguments remain open in this manuscript.

A bank is a representation unit containing $N$ weighted inverse terms. Its count is neither an arithmetic lower bound nor an evaluation cost. The generic $4N$ construction has $4N^2$ phase weights and a conservative $O(N^3)$ direct preparation/application with $O(N^2)$ stored data. Possible sharing, algebraic parameter representation, precision, stability, and weight preparation require separate analysis. No full FFT speedup, industry consequence, or historical-priority claim is inferred.

\section{Certificates and reproducibility}\label{sec:certificates}

The human-readable proof above specifies every coefficient needed for its short identities and every amplitude needed for the constructions. A companion \texttt{SOURCE\_REGISTER.json} gives full relative paths, byte lengths, and SHA-256 digests of the frozen mathematics, literal witness data, independent receipts, verifier files, and recorded outputs used in this synthesis. \texttt{THEOREM\_STATUS.json} separates accepted source results from the review status of this newly assembled manuscript. Hashes are identifiers of bytes, not proofs of the mathematics.

\paragraph{Exact data conventions.}
Cyclotomic elements use ascending coefficients of $1,q,\ldots,q^{\varphi(N)-1}$, with rational entries, reduced modulo $\Phi_N$. Polynomial lists in an auxiliary variable are also ascending. In the order-seven lower certificate, each auxiliary coefficient is a six-entry cyclotomic vector; its field named $F$ is the polynomial called $E$ in~\eqref{eq:n7F}. The short modular arrays in this manuscript are ordinary residues in $\{0,\ldots,28\}$. The order-seven upper data represent triple amplitudes by degree-at-most-two polynomials in the cubic algebra, to be evaluated at each distinct root. They do not choose approximate embeddings.

\paragraph{Order three.}
The accepted signed-magnetic packet contains the general model, support proof, construction, and an independent control that verifies all twelve literal branch matrices over several fields, including characteristic two. The all-field claim comes from the proof, not extrapolation from those examples.

\paragraph{Order five lower bound.}
The exact checker reconstructs the Gaussian numerator identity in symbolic parameters, phase orthogonality, all 25 trace coefficients, forced amplitudes, row/axis equations, type splitting, and the four nonzero elimination residues. Independent review accepts the characteristic-zero exclusion and all degeneracy coverage.

\paragraph{Order five upper bound.}
The literal file \texttt{WITNESS.json} specifies all five parameter pairs, guards, normalized amplitudes, and expanded raw weights. Its independent literal control inverts the original branches by exact elimination, verifies both inverse products and every original Fourier entry, and checks all primitive-root conjugates. A separate control verifies the nonsingular parameterization and all determinant identities.

\paragraph{Order seven lower bound.}
The accepted review's \nolinkurl{EXACT_ROW_CERTIFICATE.json} records independently derived moments, $R,E$, and the rational Bezout cofactors; \nolinkurl{LITERAL_MODEL_RESULTS.json} records direct model and symmetry checks. The frozen author data supply the modular quotient matrix and inverse, both matched independently. The independent controls separately verify the exact identity and the good-reduction certificate, including surviving denominators and leading coefficients. The proof applies to arbitrary characteristic-zero extension nodes.

\paragraph{Order seven upper bound.}
The literal \texttt{EXACT\_WITNESS.json} includes the cubic, its discriminant, each transformed amplitude vector, the reduced sector polynomials, inverse guards in the cubic algebra, and exact determinant/entry checks. The independent literal trace control derives its own compatible amplitudes. Distinct free-coordinate choices can produce different amplitude vectors while defining the same legal target; source-data comparison is therefore done at the level of the documented witness and claims, not by requiring unrelated choices to coincide.

\paragraph{Reproduction boundary.}
This synthesis has not executed, imported, compiled, or parsed any source research program. It reads mathematics, literal JSON data, receipts, and recorded outputs, and verifies source-byte hashes. Mathematical reproduction is possible by implementing the displayed field arithmetic and identities independently. For an authorized reproduction using the archived verifier programs, use a new writable copy, verify the hashes first, and run the independent controls listed in \texttt{REPRODUCIBILITY.txt}; retain all output and compare exact results. Do not run a generator in a frozen packet, because several write their result files beside themselves. Python and, where stated, SymPy are the archived control dependencies. The recipes in Sections~\ref{sec:n3}, \ref{sec:n5upper}, and \ref{sec:n7upper} do not depend on trusting a program's search history.

\section{Relation to prior work}

The underlying tools are classical. Gaussian-binomial expansion and the root-of-unity power identity precede this synthesis \cite{HMS,Schlosser}. The cyclic magnetic system of Ouvry, Wagner, and Wu \cite[Section 2.2, equations 7--8]{Ouvry} contains, after setting their reverse-shift and reverse-modulation variables to zero, the same branch matrix and denominator $1-\alpha^N-\beta^N$. This is a specialization of their system; it does not attribute the present additive bank-minimization problem to them. Finite Weyl/character algebra \cite{Gross} and Prony recurrences \cite{PottsTasche,PeterPlonka} explain the projection and rank steps.

The grammar-specific content is the isolation of coefficients by the four signed triangles, the axis coupling that excludes all order-five four-bank candidates, the two-row polynomial incompatibility at order seven with full degenerate-case coverage, and the explicit small-order legal realizations. The existing bounded prior-art comparison inspected the cited primary passages and found no equivalent exact classification in those passages. That bounded result is not evidence of worldwide priority. No new literature search was undertaken for this manuscript, and no priority claim is made.

\input{receipt_appendix.tex}

\begin{thebibliography}{9}
\bibitem{HMS} O. Holtz, V. Mehrmann, and H. Schneider,
\emph{Potter, Wielandt, and Drazin on the matrix equation $AB=\omega BA$, with some new answers to old questions},
American Mathematical Monthly \textbf{111} (2004), 655--667. Theorems 1--2.
\url{https://people.math.wisc.edu/hans/hms_sub.pdf}.
\bibitem{Schlosser} M. J. Schlosser,
\emph{A noncommutative weight-dependent generalization of the binomial theorem},
arXiv:1106.2112, version 2 (2012), Proposition 1.
\url{https://arxiv.org/pdf/1106.2112}.
\bibitem{Gross} D. Gross,
\emph{Hudson's theorem for finite-dimensional quantum systems},
Journal of Mathematical Physics \textbf{47} (2006), 122107, Section II A--D.
\url{https://arxiv.org/pdf/quant-ph/0602001}.
\bibitem{Ouvry} S. Ouvry, S. Wagner, and S. Wu,
\emph{On the algebraic area of lattice walks and the Hofstadter model},
arXiv:1606.03273 (2016), Section 2.2.
\url{https://arxiv.org/pdf/1606.03273}.
\bibitem{PottsTasche} D. Potts and M. Tasche,
\emph{Parameter estimation for exponential sums by approximate Prony method},
Signal Processing \textbf{90} (2010), 1631--1642, Section 2.
\url{https://www-user.tu-chemnitz.de/~potts/paper/prony.pdf}.
\bibitem{PeterPlonka} T. Peter and G. Plonka,
\emph{A generalized Prony method for reconstruction of sparse sums of eigenfunctions of linear operators},
author manuscript dated 10 January 2013, Theorem 2.1.
\url{https://num.math.uni-goettingen.de/plonka/pdfs/Operator_Prony_revised2.pdf}.
\end{thebibliography}
\end{document}
